% from ega0-21-fr.tex
\oldpage[0\textsubscript{IV-1}]{155}
Posant $x^{(p)}=x\otimes1$, on a donc $(\alpha x)^{(p)}=\alpha^p x^{(p)}$.
Lorsque $F_A$ est injectif, on peut identifier $A^{(p)}$ à l'anneau $A$
considéré comme algèbre sur son sous-anneau $A^p$.

\begin{proposition}[21.1.5]
\label{0.21.1.5-fr}
Soient $A$, $B$ deux anneaux, $u:A\to B$ un homomorphisme.
\begin{enumerate}
  \item[(i)] Toute $A$-dérivation de $B$ dans un $B$-module $L$ est nulle
  dans $A[B^p]$ (donc est une $A[B^p]$-dérivation).
  \item[(ii)] Pour toute sous-$A$-algèbre $A'$ de $A[B^p]$, si
  $j:A'\to B$ est l'injection canonique, l'homomorphisme canonique
  \[
  j_{B/A'/A}:\Omega^1_{B/A}\longrightarrow\Omega^1_{B/A'}
  \]
  est bijectif.
  \item[(iii)] Supposons qu'il existe un entier $s\geq0$ tel que
  $u(A)\supset B^{p^s}$. Alors, dans l'anneau $B\otimes_A B$,
  $\mathfrak J_{B/A}$ est un nilidéal.
\end{enumerate}
\end{proposition}

\begin{proof}
(i) Par récurrence à partir de (20.1.1.1), on déduit, pour toute
$A$-dérivation $D:B\to L$, que $D(x^n)=nx^{n-1}D(x)$, d'où en particulier
$D(x^p)=0$.

(ii) Vu (20.4.8), l'assertion (ii) n'est qu'une traduction de (i), cette
dernière s'écrivant $\operatorname{Der}_{A[B^p]}(B,L)=\operatorname{Der}_A(B,L)$
pour tout $B$-module $L$.

(iii) Pour tout $x\in B$, on a
$(x\otimes1-1\otimes x)^{p^s}=x^{p^s}\otimes1-1\otimes x^{p^s}=0$,
puisque $x^{p^s}\in u(A)$. La conclusion résulte de ce que les éléments
$1\otimes x-x\otimes1$ engendrent $\mathfrak J_{B/A}$ (20.4.4).
\end{proof}

Il résulte aussitôt de (21.1.5) que, pour tout couple d'homomorphismes
d'anneaux $A\to B\to C$, on a
\begin{equation}
\tag{21.1.5.1}
\label{0.21.1.5.1-fr}
\Upsilon_{C/B/A}=\Upsilon_{C/B/A'}
\end{equation}
pour toute sous-$A$-algèbre $A'\subset A[B^p]$ de $B$.

D'autre part, (21.1.5) montre en particulier que, pour les modules de
différentielles «absolues»,
\begin{equation}
\tag{21.1.5.2}
\label{0.21.1.5.2-fr}
\Omega^1_A=\Omega^1_{A/A^p}.
\end{equation}

\begin{corollary}[21.1.6]
\label{0.21.1.6-fr}
Supposons que $B$ soit une $A$-algèbre de type fini et qu'il existe un
entier $s\geq0$ tel que $u(A)\supset B^{p^s}$. Si $\Omega^1_{B/A}=0$,
$u$ est surjectif.
\end{corollary}

\begin{proof}
